\documentclass[11pt]{article}

\usepackage[margin=1in]{geometry}
\usepackage{amsmath,amssymb,amsthm,mathtools,booktabs,enumitem,microtype,xcolor}
\usepackage[colorlinks=true,linkcolor=blue!50!black,urlcolor=blue!50!black,citecolor=blue!50!black]{hyperref}

\newtheorem{theorem}{Theorem}
\newtheorem*{lemmaA}{Lemma A}
\newtheorem*{lemmaB}{Lemma B}
\newtheorem{proposition}{Proposition}
\newtheorem{corollary}{Corollary}
\theoremstyle{definition}
\newtheorem{definition}{Definition}
\theoremstyle{remark}
\newtheorem*{remark}{Remark}

\DeclareMathOperator{\sign}{sign}
\DeclareMathOperator{\maj}{maj}

\title{Refutation of Conjecture 1.4 (\texttt{1.conj.2}) on Mixed Memories in Hopfield Networks}
\author{}
\date{November 2026}

\begin{document}
\maketitle

\begin{center}
\small
\textbf{Target Reference:} V\'eronique Gayrard, \emph{Mixed memories in Hopfield networks},
\href{https://arxiv.org/abs/2504.04879}{arXiv:2504.04879} (\texttt{Mix-ALL\_v2.tex}, revised June 2026),
Conjecture~1.4 (\TeX{} label \texttt{1.conj.2}).\\[3pt]
\textbf{Domain:} Computational Neuroscience / Statistical Mechanics of Associative Memory.\\[3pt]
\textbf{Verification:} Minimal $n=7$ dense witness $((\frac{11}{32})^3, (\frac{9}{32})^4)$ for $p\in\{3,4\}$
and Parseval lifting inequality $\sum \tilde{m}_\tau^{p-1} \le 1 < 2$ certified in Lean~4;
exact rational arithmetic verification in \texttt{verify\_hopfield.py}.
\end{center}

\medskip

\begin{abstract}
Conjecture~1.4 (\texttt{1.conj.2}) of Gayrard (\href{https://arxiv.org/abs/2504.04879}{arXiv:2504.04879}) is \textbf{false}: while the ``if'' direction is Gayrard's Theorem~\texttt{1.theo2.mix}, the ``only if'' direction fails structurally. We establish explicit, strictly tie-free refutations for:
\begin{enumerate}[leftmargin=2em,itemsep=2pt]
  \item the classical Hopfield model $F(x) = \frac{1}{2}x^2$, where the conjecture holds for $n \in \{1, 3\}$ and fails for \emph{every odd $n \ge 5$};
  \item dense Hopfield models $F(x) = \frac{1}{p}x^p$, proved rigorously for $3 \le p \le 19$ and constructively verified via exact rational arithmetic for $20 \le p \le 500$, with minimal dimension $n_0(3) = n_0(4) = 7$ and lifting to all odd $n \ge n_0(p)$; and
  \item general smooth $F$ with $F'(x) > 0$ on $(0,\infty)$, failing at $n \in \{1,3\}$ when $F'$ is not odd and at $n = 5$ even for smooth odd strictly increasing $F'$ with $m \ge 0$.
\end{enumerate}
\end{abstract}

\begin{remark}[Prior Art and Novelty]
\emph{Classical Hopfield ($p = 2$) prior art:} The $n = 5$ counterexample $(\frac{5}{8}, \frac{3}{8}, \frac{3}{8}, \frac{1}{8}, \frac{1}{8})$, the even-multiplicity obstruction (Lemma~B), and the even-block lifting to all odd $n \ge 5$ were posted publicly in October 2026 (\texttt{the-omega-institute/trureturing\#13571}, PR~\texttt{\#13650}, and \texttt{\#17309}). In addition, Francesco Tosello announced a UBC seminar (November 18, 2026, \emph{Stable Mixed Memories in the Hopfield Model}, preprint not yet public) addressing stable mixed memories and Gayrard's conjecture; that work should be consulted once publicly posted.

\emph{Novel contributions here:} (i)~the dense Hopfield ($p \ge 3$) refutation---the minimal $n = 7$ witness $((\frac{11}{32})^3, (\frac{9}{32})^4)$ for $p = 3, 4$ (certified in Lean~4), the general two-block perturbed-majority family $(2a+1, 2b)$ with exact Chow parameter formulas $(\Phi_1, \Phi_2)$, the exact rational proof for $3 \le p \le 19$ and exact rational verification for $20 \le p \le 500$, and the Parseval lifting inequality $\sum \tilde{m}_\tau^{p-1} \le 1 < 2$ (also certified in Lean~4); and (ii)~the general smooth $F$ counterexamples ($F(x) = e^x - x$ at $n = 1, 3$ and $F(x) = \sqrt{1 + 4x^2}$ at $n = 5$).
\end{remark}

\section{Foundational Characterization and Multiplicity Invariant}

\begin{lemmaA}[Exact Characterization of Tie-Free $n$-Mixed Memories]
Let $F : \mathbb{R} \to \mathbb{R}$ be smooth with $F'(x) > 0$ for $x > 0$, let $n \in \mathbb{N}$ be independent of $N$, and let $m \in [-1,1]^M$ be deterministic with exactly $n$ non-zero entries $\hat{m} = (m_{\mu_1}, \dots, m_{\mu_n}) \in ([-1,1]\setminus\{0\})^n$. Set
\[
  w = \bigl(F'(\hat{m}_1), \dots, F'(\hat{m}_n)\bigr) \in \mathbb{R}^n.
\]
Suppose:
\begin{enumerate}[leftmargin=2em,itemsep=2pt]
  \item \emph{(Tie-freeness)} $(r, w) := \sum_{\nu=1}^n r_\nu w_\nu \neq 0$ for every $r \in \{-1,1\}^n$; and
  \item \emph{(Chow fixed-point equation)}
  \begin{equation}\label{eq:chow}
    \hat{m}_\nu = \Phi_\nu(w) := \frac{1}{2^n} \sum_{r \in \{-1,1\}^n} r_\nu \sign\bigl((r, w)\bigr)
    \qquad \forall\, 1 \le \nu \le n. \tag{$\star$}
  \end{equation}
\end{enumerate}
Then:
\begin{enumerate}[label=\textup{(\alph*)},leftmargin=2em,itemsep=2pt]
  \item $\xi^{(N)}(m) \in \{-1,1\}^N = \Sigma_N$ deterministically for every $N \ge 1$, and $\xi^{(N)}(m)$ is an $n$-mixed memory of type $F$ in the sense of Definition~\texttt{1.def1.2} of Gayrard (by Kolmogorov's Strong Law of Large Numbers applied to the i.i.d.\ bounded $\{-1,1\}$ random variables $\xi_i^\mu \xi_i(m)$).
  \item Letting $C_F(m) := \min_{r \in \{-1,1\}^n} |(r, w)| > 0$, Gayrard's Section~4 Borel--Cantelli proofs (Lemmas~\texttt{4.lem1}--\texttt{4.lem4bis}) depend only on~\eqref{eq:chow} and $C_F(m) > 0$. Consequently, $\xi^{(N)}(m)$ is an asymptotically almost sure local minimum (fixed point of $T^G$ and $T^{HK}$) whenever
  \[
    M(N) \le \frac{C_F(m)^2 N}{2(2+\varepsilon)\ln N} \quad (p=2),
    \qquad\text{or}\qquad
    M(N) \le \frac{C_F(m)^2 N^{p-1}}{2(2+\varepsilon)[(2p-3)!!]\ln N} \quad (p \ge 3).
  \]
\end{enumerate}
\end{lemmaA}

\begin{lemmaB}[Multiplicity Invariant of $\mathcal{M}_n^{\mathrm{all}}$]
Let $n \in \mathbb{N}$ be odd and let $m' \in \mathcal{M}_n^{\mathrm{all}}$ (in particular, any $m' \in \mathcal{M}_{n,F}$). Let $v_1 > v_2 > \dots > v_r > 0$ be the distinct positive magnitudes appearing in $|m'|$, with multiplicities $c_1, \dots, c_r$. Then
\[
  c_1, c_2, \dots, c_{r-1} \text{ are all even, and } c_r \text{ is odd.}
\]
\end{lemmaB}

\begin{proof}
By definition, $|m'|$ takes the values $\gamma^{(k)} = \prod_{l=1}^k \alpha^{(n_l)}$ with multiplicity $n_k$ ($1 \le k \le \ell$), where $n_1, \dots, n_{\ell-1} \ge 2$ are even and $n_\ell \ge 1$ is odd. Since $\alpha^{(2j)} = 2^{-2j}\binom{2j}{j} \le \frac{1}{2} < 1$ for all $j \ge 1$, we have $\gamma^{(1)} > \dots > \gamma^{(\ell-1)} > 0$. If $n_\ell = 1$, then $\alpha^{(1)} = 1$ merges $\gamma^{(\ell)} = \gamma^{(\ell-1)}$, giving $(c_1,\dots,c_r) = (n_1,\dots,n_{\ell-2}, n_{\ell-1}+1)$ (or $(1)$ if $\ell=1$). If $n_\ell \ge 3$, then $\alpha^{(n_\ell)} = \alpha^{(n_\ell-1)} \le \frac{1}{2} < 1$, giving $(c_1,\dots,c_r) = (n_1,\dots,n_{\ell-1}, n_\ell)$. In both cases $c_1,\dots,c_{r-1}$ are even and $c_r$ is odd.
\end{proof}

\section{Refutation for Classical Hopfield \texorpdfstring{$F(x) = \frac{1}{2}x^2$}{F(x) = x\textasciicircum 2/2} (Every Odd \texorpdfstring{$n \ge 5$}{n >= 5})}

\begin{theorem}\label{thm:classical}
Let $F(x) = \frac{1}{2}x^2$ ($F'(x) = x$).
\begin{enumerate}[leftmargin=2em,itemsep=2pt]
  \item For $n \in \{1, 3\}$, every tie-free $n$-mixed memory of type $F$ satisfies $m \in \mathcal{M}_{n,F}$.
  \item For every integer $k \ge 0$, setting $n = 2k + 5$ (all odd $n \ge 5$) and $M \ge n$, the vector
  \[
    m = \alpha^{(2k)}\!\left(\underbrace{1,\dots,1}_{2k},\;\frac{5}{8},\;\frac{3}{8},\;\frac{3}{8},\;\frac{1}{8},\;\frac{1}{8},\;0,\dots,0\right) \in [-1,1]^M
    \qquad \bigl(\alpha^{(0)} := 1\bigr)
  \]
  is strictly tie-free with margin $C(m) = \frac{1}{8}\alpha^{(2k)} > 0$, defines a genuine $n$-mixed memory $\xi^{(N)}(m)$ of type $F$ (and an asymptotically almost sure local minimum of $E_{N,M}$), and satisfies $m \notin \mathcal{M}_n^{\mathrm{all}}$ (hence $m \notin \mathcal{M}_{n,F}$).
\end{enumerate}
\end{theorem}

\begin{proof}
We treat the base case $k = 0$ and then lift to all $k \ge 1$.

\emph{Base case $k = 0$ ($n = 5$).} For $\tilde{m} = (\frac{5}{8}, \frac{3}{8}, \frac{3}{8}, \frac{1}{8}, \frac{1}{8})$, we have
\[
  8(r, \tilde{m}) = 5r_1 + 3r_2 + 3r_3 + r_4 + r_5 \equiv 1 \pmod{2},
\]
so $|(r, \tilde{m})| \ge \frac{1}{8} > 0$ for all $r \in \{-1,1\}^5$. Among the $16$ vectors with $r_1 = +1$, the inequality $8(r, \tilde{m}) < 0$ holds exclusively on the three vectors
\[
  B = \bigl\{(+1,-1,-1,+1,-1),\; (+1,-1,-1,-1,+1),\; (+1,-1,-1,-1,-1)\bigr\}.
\]
By oddness of $\sign$,
\begin{align*}
  \Phi_1(\tilde{m}) &= \frac{13 - 3}{16} = \frac{5}{8} = \tilde{m}_1, \\
  \Phi_2(\tilde{m}) = \Phi_3(\tilde{m}) &= \frac{0 - 2(-3)}{16} = \frac{3}{8} = \tilde{m}_2 = \tilde{m}_3, \\
  \Phi_4(\tilde{m}) = \Phi_5(\tilde{m}) &= \frac{0 - 2(-1)}{16} = \frac{1}{8} = \tilde{m}_4 = \tilde{m}_5.
\end{align*}
Its strictly decreasing positive values $\frac{5}{8} > \frac{3}{8} > \frac{1}{8}$ have multiplicities $(1, 2, 2)$, violating Lemma~B; hence $\tilde{m} \notin \mathcal{M}_5^{\mathrm{all}}$. (Its limiting energy is $-\frac{45}{128}$, strictly below Gayrard's $(2,3)$ memory at $-\frac{44}{128}$ and equal to the symmetric $(\frac{3}{8})^5$ memory.)

\emph{General case $k \ge 1$ ($n = 2k+5$).} Write $r = (r', r'') \in \{-1,1\}^{2k} \times \{-1,1\}^5$ and $S_1(r') = \sum_{\nu=1}^{2k} r'_\nu$. Since $|(r'', \tilde{m})| \le \frac{13}{8} < 2$, whenever $S_1(r') \neq 0$ we have $|S_1(r')| \ge 2$ and $\sign((r, \hat{m})) = \sign(S_1(r'))$ with margin at least $\frac{3}{8}\alpha^{(2k)}$; when $S_1(r') = 0$, we have $\sign((r, \hat{m})) = \sign((r'', \tilde{m}))$ with margin at least $\frac{1}{8}\alpha^{(2k)}$. Splitting $\sum_{r', r''}$ into $\{S_1 \neq 0\}$ and $\{S_1 = 0\}$ yields
\[
  \Phi_\nu(\hat{m}) = \alpha^{(2k)} \quad (1 \le \nu \le 2k),
  \qquad
  \Phi_{2k+\tau}(\hat{m}) = 2^{-2k}\binom{2k}{k}\tilde{m}_\tau = \alpha^{(2k)}\tilde{m}_\tau \quad (1 \le \tau \le 5).
\]
By Lemma~A, $m$ defines a genuine $n$-mixed memory and asymptotically a.s.\ local minimum, while its multiplicities $(2k, 1, 2, 2)$ violate Lemma~B.
\end{proof}

\section{Refutation for Dense Hopfield Models \texorpdfstring{$F(x) = \frac{1}{p}x^p$}{F(x) = x\textasciicircum p/p} (\texorpdfstring{$p \ge 3$}{p >= 3})}

\begin{theorem}\label{thm:dense}
For every integer $p \in [3, 500]$ (proved analytically/symbolically for $3 \le p \le 19$ via six explicit $(a, b)$ pairs, and verified via exact rational arithmetic for $20 \le p \le 500$; numerically observed to cover all $p \lesssim 8.8 \times 10^{13}$), there exists an odd $n_0(p) \ge 7$ such that for every odd $n \ge n_0(p)$ there is a strictly tie-free $n$-mixed memory $m \in [0,1]^M$ of type $F(x) = \frac{1}{p}x^p$ (an asymptotically almost sure local minimum for $T^G$ and $T^{HK}$) with $m \notin \mathcal{M}_n^{\mathrm{all}}$. For $p \in \{3, 4\}$, $n_0(p) = 7$ via
\[
  \tilde{m} = \left(\frac{11}{32},\;\frac{11}{32},\;\frac{11}{32},\;\frac{9}{32},\;\frac{9}{32},\;\frac{9}{32},\;\frac{9}{32}\right) \notin \mathcal{M}_7^{\mathrm{all}}.
\]
\end{theorem}

\begin{proof}[Proof Summary]
The construction proceeds in three steps.
\begin{enumerate}[leftmargin=2em,itemsep=4pt]
  \item \emph{Two-block perturbed-majority family $(2a+1, 2b)$ ($a \ge 1$, $b \ge a+1$, $n_0 = 2a+2b+1$).}
  Consider weights $w = (A^{2a+1}, B^{2b})$ satisfying
  \begin{equation}\label{eq:ratio-band}
    \frac{2a+2}{2a+1} < \frac{A}{B} < \frac{2a}{2a-1}. \tag{R}
  \end{equation}
  Under~\eqref{eq:ratio-band}, $(r, w) \neq 0$ for all $r \in \{-1,1\}^{n_0}$, and $\sign((r,w)) \neq \maj(r)$ if and only if $r \in E_+ \cup (-E_+)$, where
  \[
    E_+ = \bigl\{r = (r', r'') : S_1(r') = 2a+1,\; S_2(r'') = -(2a+2)\bigr\},
    \qquad |E_+| = \binom{2b}{b-a-1}.
  \]
  Hence the exact Chow parameters are $\Phi(w) = (\Phi_1^{2a+1}, \Phi_2^{2b})$ with
  \begin{align*}
    \Phi_1(a,b) &= \alpha^{(n_0)}\bigl(1 + x(a,b)\bigr),
    & \Phi_2(a,b) &= \alpha^{(n_0)}\!\left(1 - \frac{a+1}{b}\,x(a,b)\right), \\
    x(a,b) &= \frac{2\binom{2b}{b-a-1}}{\binom{2a+2b}{a+b}}.
  \end{align*}
  Setting $q(a,b) := \Phi_1(a,b)/\Phi_2(a,b)$, whenever
  \[
    \frac{2a+2}{2a+1} < q(a,b)^{p-1} < \frac{2a}{2a-1},
  \]
  the vector $\tilde{m}(a,b) = (\Phi_1^{2a+1}, \Phi_2^{2b})$ solves~\eqref{eq:chow} for $F(x) = \frac{1}{p}x^p$ and has multiplicities $(2a+1, 2b)$ (odd, even), violating Lemma~B.

  \item \emph{Universal Parseval lifting ($n_0 \mapsto n_0 + 2k$ for all $k \ge 1$).}
  For any tie-free solution $\tilde{m} \in (0,1)^{n_0}$ of~\eqref{eq:chow} and any $p \ge 3$, Parseval's identity for the Boolean function $g(r'') = \sign((r'', \tilde{m}^{p-1}))$ yields
  \[
    \sum_{\tau=1}^{n_0} \tilde{m}_\tau^{p-1} \le \sum_{\tau=1}^{n_0} \tilde{m}_\tau^2 \le 1 < 2.
  \]
  Therefore $\hat{m} = \alpha^{(2k)}(1^{2k}, \tilde{m})$ is strictly tie-free, solves~\eqref{eq:chow} for $F(x) = \frac{1}{p}x^p$, and violates Lemma~B (having multiplicity sequence $(2k, 2a+1, 2b)$).

  \item \emph{Explicit coverage in $p$.}
  \begin{itemize}[leftmargin=1.5em,itemsep=2pt]
    \item \emph{$3 \le p \le 19$ (proved via exact finite rational checks):} Covered explicitly by the six $(a,b)$ pairs listed in Table~\ref{tab:dense-pairs}.
    \item \emph{$20 \le p \le 500$ (exact rational verification) and higher $p$ (numerical):} The script \texttt{verify\_\allowbreak hopfield.py} constructs and verifies an exact rational pair $(a, b)$ satisfying~\eqref{eq:ratio-band} for every integer $p \in [20, 500]$. Numerical monotonicity and interval-overlap checks of $q(a, b)$ cover $a = 4,\dots,25$ up to $p \approx 8.8 \times 10^{13}$ (an explicit analytic proof of the monotonicity and interval-overlap inequalities for all $a \ge 6$ is not yet written out).
  \end{itemize}
\end{enumerate}
\end{proof}

\begin{table}[ht]
\centering
\caption{Explicit $(a, b)$ pairs covering $3 \le p \le 19$ in Theorem~\ref{thm:dense}, with base dimension $n_0 = 2a + 2b + 1$ and exact Chow ratio $q(a,b) = \Phi_1(a,b)/\Phi_2(a,b)$.}
\label{tab:dense-pairs}
\begin{tabular}{@{}cccc@{}}
\toprule
$(a, b)$ & $n_0 = 2a+2b+1$ & Exact ratio $q(a,b)$ & Ratio band $\bigl(\frac{2a+2}{2a+1}, \frac{2a}{2a-1}\bigr)$ \\
\midrule
$(1, 2)$ & $7$  & $11/9$            & $(4/3,\; 2)$ \\
$(2, 6)$ & $17$ & $121/115$         & $(6/5,\; 4/3)$ \\
$(2, 5)$ & $15$ & $587/563$         & $(6/5,\; 4/3)$ \\
$(2, 4)$ & $13$ & $235/228$         & $(6/5,\; 4/3)$ \\
$(3, 13)$ & $33$ & $138977/137107$  & $(8/7,\; 6/5)$ \\
$(2, 3)$ & $11$ & $127/125$         & $(6/5,\; 4/3)$ \\
\bottomrule
\end{tabular}
\end{table}

\section{General Smooth \texorpdfstring{$F$}{F} Counterexamples and Conceptual Diagnosis}

\begin{enumerate}[leftmargin=2em,itemsep=4pt]
  \item \emph{Sign-parity defect at $n = 1, 3$.} For $F(x) = e^x - x$ ($F'(x) = e^x - 1 > 0$ on $(0,\infty)$, not odd), Eq.~\texttt{(1.2.1.15)} of Gayrard restricts $\mathcal{M}_{n,F}$ to non-negative vectors, yet
  \[
    m = (-1, 0, \dots, 0) \quad (n=1)
    \qquad\text{and}\qquad
    m = \Bigl(-\frac{1}{2}, -\frac{1}{2}, -\frac{1}{2}, 0, \dots, 0\Bigr) \quad (n=3)
  \]
  are genuine tie-free mixed memories of type $F$.

  \item \emph{Odd strictly increasing $F'$, $m \ge 0$, $n = 5$.} For $F(x) = \sqrt{1 + 4x^2}$, the vector
  \[
    \hat{m} = \left(\frac{7}{8},\;\frac{1}{8},\;\frac{1}{8},\;\frac{1}{8},\;\frac{1}{8}\right) \notin \mathcal{M}_5^{\mathrm{all}}
  \]
  satisfies $2 < F'(7/8)/F'(1/8) = \sqrt{833/65} < 4$ and solves~\eqref{eq:chow} with zero ties.

  \item \emph{Conceptual diagnosis of the failure of \texttt{1.conj.2}.} Tie-free $n$-mixed memories are fixed points $\hat{m} = \Phi(F'(\hat{m}))$ of the Chow-parameter map on odd linear threshold functions (LTFs). Gayrard's set $\mathcal{M}_n^{\mathrm{all}}$ enumerates only \emph{lexicographic} (hierarchical-majority) LTFs. All non-degenerate odd LTFs are lexicographic for $n = 1, 3$, but non-lexicographic odd LTFs appear at $n = 5$ (such as $(3,2,2,1,1)$ for $p=2$ and $(3,1,1,1,1)$) and at $n = 7$ (such as $(11^3, 9^4)$ for $p = 3, 4$), persisting for every polynomial $F(x) = x^p/p$.
\end{enumerate}

\end{document}
